\section{Review and correction of the claimed ``disproof'' for Erd\H{o}s Problem \#65}

\subsection*{Verdict}

The arithmetic observation in the proposed write-up is correct, but it does \emph{not} solve Erd\H{o}s Problem \#65 as stated in the original source. It only disproves an \emph{over-generalised reformulation} that asks for a complete bipartite minimiser for \emph{every} exact pair $(n,m)$. The original 1981 strengthening of Erd\H{o}s is narrower: it asks about graphs with exactly $k(n-k)$ edges, i.e. precisely the edge counts realised by $K_{k,n-k}$.

Accordingly, the claims
\begin{quote}
\textbf{``FULL SOLUTION --- COUNTEREXAMPLE / DISPROOF''}
\qquad\text{and}\qquad
\textbf{``COMPLETION: 100\%''}
\end{quote}
are incorrect for the original open problem.

\subsection*{1. Three different formulations that must be separated}

The literature contains three related but distinct questions.

\begin{enumerate}
\item[(A)] \textbf{The 1975 exact-edge minimisation problem.}
Erd\H{o}s wrote: for graphs $G(n;\lfloor kn\rfloor)$, determine or estimate the minimum of
\[
\sum_{r\in \mathcal C(G)} \frac1r.
\]
He said it seemed likely that this minimum is $(\tfrac12+o(1))\log k$.

\item[(B)] \textbf{The 1981 strengthened conjecture.}
Erd\H{o}s then proposed a sharper statement: if $n\ge 2k$, then for every graph $G(n;k(n-k))$,
\[
\sum_{u_i\in \mathcal C(G)} \frac1{u_i}
\ge \frac12\left(\frac12+\frac13+\cdots+\frac1k\right),
\]
and in his own words this means the sum is \emph{minimum for the complete bipartite graph} with parts of sizes $k$ and $n-k$.

\item[(C)] \textbf{The over-generalised exact-$(n,m)$ reformulation.}
For every fixed exact pair $(n,m)$, minimise $S(G)$ over all $n$-vertex $m$-edge graphs and ask whether the minimiser is complete bipartite.
\end{enumerate}

The proposed write-up correctly disproves (C), but the original conjectural minimiser statement in Erd\H{o}s 1981 is (B), not (C).

\subsection*{2. What is correct in the proposed write-up}

The following observation is valid.

\textbf{Proposition 2.1.}
Let $n=2t+1$ with $t\ge 2$. A complete bipartite graph on $n$ vertices has exactly $a(n-a)$ edges for some integer $a$, so the possible edge counts are
\[
a(2t+1-a)\qquad(0\le a\le 2t+1).
\]
The largest such value is $t(t+1)$, attained at $a=t,t+1$, and the next largest is
\[
(t-1)(t+2)=t(t+1)-2.
\]
Hence there is \emph{no} complete bipartite graph on $2t+1$ vertices with exactly
\[
m=t(t+1)-1
\]
edges.

Therefore, if one asks the universal exact-edge question (C), then that question is false: the class of graphs with $(n,m)=(2t+1,t(t+1)-1)$ is nonempty and finite, so a minimiser of $S(G)$ exists, but no complete bipartite graph lies in that class.

\subsection*{3. Why this does not solve the original open problem}

The counterexample above does \emph{not} apply to Erd\H{o}s's 1981 strengthening, because that conjecture only concerns edge counts of the form
\[
m=k(n-k),
\]
exactly the counts realised by the candidate extremal graph $K_{k,n-k}$.

So the right conclusion is:
\begin{quote}
The proposed argument refutes only the over-broad formulation (C). It does not refute, and does not settle, the original conjecture (B).
\end{quote}

In particular, the sentence
\begin{quote}
\emph{``The minimal corrected exact-edge question is then: whenever $m=a(n-a)$, is the minimum attained by $K_{a,n-a}$?''}
\end{quote}
is not a new correction at all: it is essentially Erd\H{o}s's original 1981 strengthened conjecture.

\subsection*{4. What the recent survey actually gives}

Montgomery's 2025 survey states that in forthcoming work with Milojevi\'c, Pokrovskiy, and Sudakov, there is a threshold $d_0$ such that for every $d\ge d_0$, among all $n$-vertex graphs with at least $d(n-d)$ edges,
\[
\sum_{\ell\in \mathcal C(G)} \frac1\ell
\]
is minimised exactly by the complete bipartite graph $K_{d,n-d}$.

This is a \emph{threshold} theorem, but it also immediately implies the corresponding \emph{exact-edge} statement for the special edge count $m=d(n-d)$: since $K_{d,n-d}$ belongs to the class of graphs with exactly $d(n-d)$ edges, uniqueness in the larger ``at least'' class forces it also to be the minimiser in the exact class.

Thus the survey gives genuine partial progress on the original 1981 conjecture, for all sufficiently large $d$.

\subsection*{5. Specific mistakes in the original draft}

\begin{enumerate}
\item The phrase ``under the exact-edge interpretation of Problem~65'' is misleading. The original 1981 source does \emph{not} formulate the complete-bipartite conjecture for arbitrary exact pairs $(n,m)$.

\item Theorem~5.3 is mathematically correct only as a disproof of the over-generalised formulation (C), not of Erd\H{o}s Problem~65 in its original strengthened form.

\item The sentence calling the ``corrected exact-edge question'' the version with $m=a(n-a)$ is historically inaccurate: that is already the original 1981 conjecture.

\item The note says the large-$d$ survey theorem only addresses a threshold formulation. That is incomplete. It also yields the exact-edge statement for $m=d(n-d)$.

\item The final verdict ``FULL SOLUTION --- COUNTEREXAMPLE / DISPROOF'' is false as a verdict on the original open problem.

\item The completion estimate ``100\%'' is false for the original problem.

\item The bibliographic reference ``EMS Magazine \textbf{138} (2025/2026 issue)'' should be corrected to \emph{EMS Magazine} \textbf{138} (2025).

\item If one keeps a statement ranging over $0\le a\le n$, then degenerate cases $a=0,n$ should be separated from formulas involving harmonic numbers. For cycle questions, the natural nontrivial range is $1\le a\le n-1$.
\end{enumerate}

\subsection*{6. Corrected final conclusion}

\textbf{Correct conclusion.}
The proposed note does \emph{not} provide a satisfactory solution to Erd\H{o}s Problem \#65. What it does show is the following:
\begin{quote}
If one over-generalises the minimiser question to arbitrary exact edge classes $(n,m)$, then that over-generalised statement is trivially false because many edge counts are not realised by any complete bipartite graph on $n$ vertices.
\end{quote}
This is a valid clarification, but it is not a solution of the original conjecture, which already restricts to the exact edge count $m=k(n-k)$.

The best current picture from the checked sources is therefore:
\begin{itemize}
\item the broad lower-bound problem was introduced earlier and has asymptotically sharp lower bounds;
\item Erd\H{o}s's 1981 complete-bipartite strengthening remains open in general;
\item Montgomery's 2025 survey reports a forthcoming exact extremal theorem proving the conjectured complete-bipartite minimiser for all sufficiently large $d$, in the class of graphs with at least $d(n-d)$ edges, and hence also for the exact edge count $d(n-d)$.
\end{itemize}

\subsection*{7. References}

\begin{thebibliography}{9}

\bibitem{Erdos1975}
P.~Erd\H{o}s,
\emph{Problems and results on finite and infinite graphs},
available at the R\'enyi archive (1975).

\bibitem{Erdos1981}
P.~Erd\H{o}s,
\emph{On the combinatorial problems which I would most like to see solved},
Combinatorica \textbf{1} (1981), 25--42.

\bibitem{GPVS1985}
A.~Gy\'arf\'as, H.~J.~Pr\"omel, B.~Voigt, and E.~Szemer\'edi,
\emph{On the sum of the reciprocals of cycle lengths in sparse graphs},
Combinatorica \textbf{5} (1985), 41--52.

\bibitem{Montgomery2025}
R.~Montgomery,
\emph{Cycles and expansion in graphs},
EMS Magazine \textbf{138} (2025), 5--12.

\bibitem{Problem65Page}
T.~F.~Bloom,
\emph{Erd\H{o}s Problem \#65}, online problem page, accessed 2026-03-09.

\end{thebibliography}
